\section{System overview}
\label{sec:system-overview}

\subsection{Architecture}

Meowtion is a four-tier system: the collar on the cat, the station in the home, the cloud
backend, and the owner's dashboard. \Cref{fig:architecture} shows the data path.

\begin{figure}[ht]
  \centering
  \begin{tikzpicture}[node distance=8mm and 12mm]
    \node[mwnode] (cat) {Cat};
    \node[mwnode, right=of cat, fill=mwlav!10, draw=mwlav] (collar)
      {\textbf{Collar}\\\footnotesize XIAO nRF52840 Sense\\\footnotesize Zephyr, battery\\\footnotesize on-device AI};
    \node[mwnode, right=of collar, fill=mwlav!10, draw=mwlav] (station)
      {\textbf{Station}\\\footnotesize XIAO ESP32-S3\\\footnotesize ESP-IDF, mains\\\footnotesize Wi-Fi gateway};
    \node[mwnode, right=of station] (cloud)
      {\textbf{Firebase}\\\footnotesize Auth, Database\\\footnotesize Storage, Functions};
    \node[mwnode, below=14mm of station] (dash)
      {\textbf{Dashboard}\\\footnotesize Streamlit + web SDK};
    \node[mwnode, left=of dash] (owner) {Owner};

    \draw[mwedge] (cat) -- (collar);
    \draw[mwedge] (collar) -- node[above,font=\scriptsize,text=mwmut]{BLE} (station);
    \draw[mwedge] (station) -- node[above,font=\scriptsize,text=mwmut]{Wi-Fi / HTTPS} (cloud);
    \draw[mwedge] (cloud.south) |- (dash.east);
    \draw[mwedge] (dash) -- (owner);
  \end{tikzpicture}
  \caption{End-to-end data path. The collar is BLE-only, so the always-on station is its
           gateway to the cloud.}
  \label{fig:architecture}
\end{figure}

\subsection{Data flow}

In normal operation the collar runs a classifier on its own sensors and advertises a compact
telemetry record over BLE (\cref{sec:protocols}). The station, scanning continuously, reads
that record, attaches the cat and device identity, and writes it to the cloud over an
authenticated HTTPS connection. The dashboard reads the same data and renders the owner's view.
Because the collar transmits only a classification result and not raw sensor data, the uplink is
tiny and the cat's audio never leaves the device.

The system has two operating modes, selected per station from the dashboard:

\begin{description}[leftmargin=1.6em,style=nextline]
  \item[Training mode] The station records short labelled audio clips (with time-aligned motion)
        while a registered collar is in range, and uploads them for model training. This is how a
        new behaviour is taught.
  \item[Production mode] The station stops recording entirely and simply relays the collar's
        on-device classification. This is the everyday monitoring mode.
\end{description}

\subsection{Multi-user model}

Meowtion is multi-tenant. Each owner creates an account, registers their own station(s) and
collar(s), and sees only their own cats. Identity and isolation are enforced by the backend's
authentication and security rules (\cref{sec:cloud-backend,sec:security-privacy}), not by the
client. A single public, read-only demo account lets anyone explore the dashboard without
sign-in.

\subsection{Technology summary}

\begin{table}[ht]
  \centering
  \caption{Technology at each tier.}
  \label{tab:stack}
  \begin{tabularx}{\textwidth}{@{}lX@{}}
    \toprule
    \textbf{Tier} & \textbf{Stack} \\
    \midrule
    Collar    & Seeed XIAO nRF52840 Sense; Zephyr / nRF Connect SDK; TensorFlow Lite Micro \\
    Station   & Seeed XIAO ESP32-S3; ESP-IDF; NimBLE (BLE central) + Wi-Fi \\
    Cloud     & Firebase: Authentication, Realtime Database, Cloud Storage, Python Cloud Functions \\
    Dashboard & Streamlit (Python) + Firebase web SDK (JavaScript) \\
    Training  & TensorFlow, executed server-side in a Cloud Function \\
    \bottomrule
  \end{tabularx}
\end{table}
